% chapter1.tex
% ============================================================
%  فصل ۱: مبانی نظری شکل‌گیری دولت
% ============================================================

\chapter{مبانی نظری شکل‌گیری دولت}

\begin{keybox}[خلاصه‌ی فصل]
در این فصل با مفاهیم بنیادین \keyword{دولت} (\en{State})، \keyword{ملت} (\en{Nation})
و \keyword{قومیت} (\en{Ethnicity}) آشنا می‌شویم و نظریه‌های اصلی شکل‌گیری دولت‌ها را
در حوزه‌های جغرافیایی مختلف مرور می‌کنیم.
\end{keybox}

% ============================================================
\section{مفهوم‌شناسی: دولت، ملت و قومیت}
% ============================================================

\subsection{دولت (\en{State})}

\lettrine[lines=2, findent=3pt, nindent=0pt]{\textcolor{PrimaryDark}{د}}{}%
ولت در معنای مدرن خود — آن‌چنان‌که \concept{ماکس وبر} تعریف کرده — نهادی است که
\important{انحصار مشروع خشونت} را در قلمرو جغرافیایی معیّنی در اختیار دارد.
اما این تعریف تنها نقطه‌ی آغاز است. چهار عنصر کلاسیک دولت عبارت‌اند از:

\begin{enumerate}
  \item \keyword{قلمرو مشخص} (\en{Territory})
  \item \keyword{جمعیت دائمی} (\en{Population})
  \item \keyword{حکومت} (\en{Government})
  \item \keyword{حاکمیت} (\en{Sovereignty})
\end{enumerate}

\begin{figure}[htbp]
\centering
\begin{tikzpicture}[
  every node/.style={font=\small},
  element/.style={
    draw=PrimaryMid, fill=BgLightBlue, rounded corners=8pt,
    minimum width=3.2cm, minimum height=1.4cm, align=center,
    font=\small\bfseries, text=PrimaryDark,
    drop shadow={shadow xshift=1.5pt, shadow yshift=-1.5pt, opacity=0.15}
  },
  center/.style={
    circle, draw=AccentGold, fill=BgLightGold, line width=2pt,
    minimum size=2.8cm, align=center, font=\bfseries\large,
    text=PrimaryDark,
    drop shadow={shadow xshift=2pt, shadow yshift=-2pt, opacity=0.2}
  },
  arrow/.style={
    -{Stealth[length=6pt]}, line width=1.2pt, PrimaryMid
  }
]
  \node[center] (state) {دولت\\\en{\footnotesize State}};

  \node[element, above=1.8cm of state] (territory)
    {قلمرو\\\en{\footnotesize Territory}};
  \node[element, below=1.8cm of state] (pop)
    {جمعیت\\\en{\footnotesize Population}};
  \node[element, left=2.2cm of state] (gov)
    {حکومت\\\en{\footnotesize Government}};
  \node[element, right=2.2cm of state] (sov)
    {حاکمیت\\\en{\footnotesize Sovereignty}};

  \draw[arrow] (state) -- (territory);
  \draw[arrow] (state) -- (pop);
  \draw[arrow] (state) -- (gov);
  \draw[arrow] (state) -- (sov);
\end{tikzpicture}
\caption{چهار عنصر سازنده‌ی دولت مدرن}
\label{fig:state-elements}
\end{figure}

\subsection{ملت (\en{Nation})}

مفهوم ملت به دو شکل عمده تعریف شده است:

\begin{keybox}[دو رویکرد به ملت]
\begin{description}[style=nextline, leftmargin=2cm]
  \item[\textcolor{PrimaryDark}{رویکرد مدنی (\en{Civic Nation})}]
    ملت بر پایه‌ی شهروندی مشترک، قانون اساسی واحد و اراده‌ی سیاسی مشترک تعریف می‌شود.
    نمونه: فرانسه، ایالات متحده.
  \item[\textcolor{AccentCopper}{رویکرد قومی (\en{Ethnic Nation})}]
    ملت بر پایه‌ی تبار، زبان، فرهنگ و تاریخ مشترک تعریف می‌شود.
    نمونه: آلمان (پیش از ۱۹۴۵)، ژاپن.
\end{description}
\end{keybox}

\subsection{قومیت (\en{Ethnicity})}

قومیت به گروهی اجتماعی اشاره دارد که اعضای آن بر پایه‌ی \concept{هویت فرهنگی مشترک} — شامل زبان، آداب، دین یا تبار مشترک — خود را از دیگران متمایز می‌دانند. قومیت لزوماً با ملت یکسان نیست. یک دولت ممکن است چندین قومیت و حتی چندین ملت را در بر بگیرد.

\simplesep

% ============================================================
\section{نظریه‌های شکل‌گیری دولت}
% ============================================================

\subsection{رویکردهای کلاسیک}

\begin{table}[htbp]
\centering
\caption{نظریه‌های اصلی شکل‌گیری دولت}
\label{tab:state-formation-theories}
\rowcolors{2}{BgLightBlue!40}{white}
\begin{tabularx}{\textwidth}{>{\bfseries\color{PrimaryDark}}r X >{\small\itshape\color{TextMid}}r}
\toprule
\rowcolor{PrimaryDark}
\textcolor{white}{\bfseries نظریه} &
\textcolor{white}{\bfseries شرح اجمالی} &
\textcolor{white}{\bfseries اندیشمند} \\
\midrule
قرارداد اجتماعی &
دولت حاصل توافق ارادی افراد برای خروج از وضع طبیعی است.
& هابز، لاک، روسو \\
فتح و غلبه &
دولت در نتیجه‌ی تسلط گروهی بر گروه دیگر شکل گرفته.
& ابن‌خلدون، اوپنهایمر \\
نظریه‌ی مارکسیستی &
دولت ابزار سلطه‌ی طبقاتی است که همراه مالکیت خصوصی پدید آمد.
& مارکس، انگلس، لنین \\
نظریه‌ی جنگ &
فشار نظامی و رقابت بین واحدها دولت‌های مدرن را ساخت.
& چارلز تیلی \\
نظریه‌ی هیدرولیکی &
مدیریت آبیاری در مقیاس بزرگ به دولت متمرکز انجامید.
& ویتفوگل \\
نهادگرایی تاریخی &
مسیرهای تاریخی ویژه (وابستگی به مسیر) ساختار دولت را شکل داد.
& برینگتن مور، اسکاچپل \\
\bottomrule
\end{tabularx}
\end{table}

\subsection{نظریه‌ی چارلز تیلی: «جنگ دولت می‌سازد»}

\begin{quotebox}
\en{``War made the state, and the state made war.''}
\hfill — \en{Charles Tilly}
\end{quotebox}

\concept{تیلی} استدلال می‌کند که رقابت نظامی قرون ۱۶ تا ۱۹ میلادی در اروپا، دولت‌ها را وادار به ایجاد بوروکراسی مالیاتی، ارتش دائمی و دستگاه اداری کارآمد کرد. دولت‌هایی که نتوانستند این سه کارکرد را توسعه دهند، از صحنه حذف شدند.

\begin{figure}[htbp]
\centering
\begin{tikzpicture}[
  node distance=1.5cm and 2.5cm,
  block/.style={
    rectangle, rounded corners=6pt, draw=PrimaryMid, fill=BgLightBlue,
    minimum height=1cm, minimum width=2.8cm, align=center,
    font=\small\bfseries, text=PrimaryDark
  },
  arrow/.style={-{Stealth[length=5pt]}, thick, PrimaryDark}
]
  \node[block] (war) {جنگ و رقابت\\نظامی};
  \node[block, right=of war] (tax) {نظام مالیاتی\\گسترده‌تر};
  \node[block, right=of tax] (bur) {بوروکراسی\\حرفه‌ای};
  \node[block, below=of bur] (army) {ارتش دائمی\\و بزرگ‌تر};
  \node[block, below=of war] (state) {دولت مدرن\\متمرکز};

  \draw[arrow] (war) -- (tax);
  \draw[arrow] (tax) -- (bur);
  \draw[arrow] (bur) -- (army);
  \draw[arrow] (army) -- (state);
  \draw[arrow, dashed, AccentCopper] (state) -- (war)
    node[midway, left, font=\footnotesize\color{AccentCopper}] {بازخورد};
\end{tikzpicture}
\caption{چرخه‌ی تیلی: جنگ و دولت‌سازی}
\label{fig:tilly-cycle}
\end{figure}

\subsection{مسیرهای مختلف دولت‌سازی در جهان}

\begin{goldbox}[سه مسیر عمده]
\begin{enumerate}
  \item \keyword{مسیر اروپایی:} جنگ → مالیات → بوروکراسی → دموکراسی تدریجی
  \item \keyword{مسیر پسااستعماری:} مرزهای مصنوعی → دولتِ وارداتی → بحران مشروعیت
  \item \keyword{مسیر انقلابی:} انقلاب → دولت ایدئولوژیک → اصلاحات یا فروپاشی
\end{enumerate}
\end{goldbox}

% ============================================================
\section{شکل‌گیری دولت‌ها در حوزه‌های جغرافیایی مختلف}
% ============================================================

\subsection{اروپا: از فئودالیسم تا دولت وستفالیایی}

صلح وستفالی (۱۶۴۸) نقطه‌ی عطفی در تاریخ شکل‌گیری نظام دولت‌های مدرن بود. با پایان جنگ‌های سی‌ساله، اصل \keyword{حاکمیت سرزمینی} (\en{Territorial Sovereignty}) به رسمیت شناخته شد. از آن پس، هر واحد سیاسی در درون مرزهای خود مرجع نهایی اقتدار بود.

\begin{figure}[htbp]
\centering
\begin{tikzpicture}[
  timenode/.style={
    draw=PrimaryMid, fill=BgLightBlue, rounded corners=4pt,
    minimum width=2.5cm, minimum height=0.8cm,
    font=\footnotesize\bfseries, text=PrimaryDark, align=center
  },
  yearnode/.style={
    fill=AccentGold, text=PrimaryDark, rounded corners=2pt,
    font=\footnotesize\bfseries, inner sep=3pt
  }
]
  % Timeline line
  \draw[PrimaryMid, line width=2pt] (0,0) -- (14,0);

  % Events
  \foreach \x/\year/\event in {
    1/1648/صلح وستفالی,
    4/1789/انقلاب فرانسه,
    7/1815/کنگره‌ی وین,
    10/1919/معاهده ورسای,
    13/1945/سازمان ملل
  }{
    \node[yearnode] at (\x, 0.5) {\year};
    \node[timenode, above=0.8cm] at (\x, 0.5) {\event};
    \draw[PrimaryMid, thick] (\x, 0) -- (\x, 0.35);
  }
\end{tikzpicture}
\caption{خط زمانی تحول نظام دولت‌ها در اروپا}
\label{fig:europe-timeline}
\end{figure}

\subsection{خاورمیانه: دولت‌سازی مصنوعی}

\begin{challengebox}[چالش مرزهای سایکس–پیکو]
مرزهای کنونی بسیاری از کشورهای خاورمیانه نه بر اساس واقعیت‌های قومی و فرهنگی، بلکه بر پایه‌ی توافق‌های استعماری (به‌ویژه سایکس–پیکو ۱۹۱۶) ترسیم شدند. نتیجه: دولت‌هایی چندقومی با هویت ملی شکننده.

نمونه‌ها: عراق (عرب‌شیعه، عرب‌سنی، کرد)، سوریه، لبنان.
\end{challengebox}

\subsection{آفریقا: میراث استعمار و دولت‌های ضعیف}

در آفریقا نیز مرزهای دوران استعمار بی‌اعتنا به واقعیت‌های قبیله‌ای و زبانی ترسیم شدند. پس از استقلال، بسیاری از دولت‌ها با بحران \concept{ملت‌سازی} مواجه شدند: چگونه از ده‌ها یا صدها قومیت، یک هویت ملّی واحد بسازیم؟

\subsection{آسیا: تنوع در مسیرهای دولت‌سازی}

\begin{itemize}
  \item \keyword{چین:} امپراتوری کهن → دولت‌ملت کمونیستی متمرکز
  \item \keyword{هندوستان:} تنوع عظیم + دموکراسی فدرال
  \item \keyword{ژاپن:} همگنی نسبی قومی + مدرن‌سازی سریع (اصلاحات میجی)
  \item \keyword{ایران:} امپراتوری کهن + تنوع قومی + تمرکزگرایی مدرن (پهلوی)
\end{itemize}

% ============================================================
\section{سرفصل‌های مرتبط در نظریه‌ی سیاسی}
% ============================================================

\begin{table}[htbp]
\centering
\caption{سرفصل‌های اصلی در حوزه‌ی دولت‌سازی و مدیریت تنوع}
\label{tab:subfields}
\rowcolors{2}{BgLightGold!50}{white}
\begin{tabularx}{\textwidth}{>{\bfseries\color{PrimaryDark}}r X}
\toprule
\rowcolor{PrimaryDark}
\textcolor{white}{\bfseries حوزه‌ی مطالعاتی} &
\textcolor{white}{\bfseries موضوعات کلیدی} \\
\midrule
سیاست تأسیسی \en{(Constitutional Design)} &
طراحی قانون اساسی، توزیع قدرت، نهادسازی \\
سیاست تطبیقی \en{(Comparative Politics)} &
مقایسه‌ی نظام‌ها، فدرالیسم تطبیقی، الگوهای دموکراسی \\
ملی‌گرایی و قومیت‌شناسی &
ملت‌سازی، ناسیونالیسم، مدیریت تنوع قومی \\
حقوق اقلیت‌ها &
حقوق زبانی، خودمختاری فرهنگی، نمایندگی \\
توسعه‌ی سیاسی &
دموکراتیزاسیون، پایداری نهادی، حکمرانی خوب \\
\bottomrule
\end{tabularx}
\end{table}

\simplesep

\begin{lessonbox}[جمع‌بندی فصل اول]
\begin{itemize}
  \item دولت‌ها در مسیرهای تاریخی بسیار متفاوتی شکل گرفته‌اند.
  \item تمایز میان دولت، ملت و قومیت برای فهم تنوع سیاسی ضروری است.
  \item بسیاری از چالش‌های امروز (عراق، سوریه، ایران) ریشه در عدم تطابق دولت و ملت دارند.
  \item نظریه‌های شکل‌گیری دولت ابزارهای تحلیلی قدرتمندی برای فهم وضعیت کنونی فراهم می‌کنند.
\end{itemize}
\end{lessonbox}